\chapter{Introduction}
\label{ch:intro}

\section{What the admin bot does}

The department runs its own chat server, \server, based on the open \emph{Matrix} protocol.
You use it with the app \emph{Element} (web, desktop or phone). A course on the server is a
\emph{space}: a folder of chat rooms that belong together, such as \enquote{Announcements},
\enquote{Questions} and one room per exercise group.

Setting this up by hand means creating rooms, configuring who may post where, and inviting
every student one by one. The admin bot \botid{} does this for you. You send it short text
commands in a direct chat, and it

\begin{itemize}
  \item creates a course space with all its rooms in one command (\autoref{ch:courses}),
  \item appoints your teaching assistants (TAs) and gives them the right permissions,
  \item brings students in: existing accounts get an invitation, new people get a personal
        signup link that you send them (\autoref{ch:people}),
  \item creates exercise group rooms, fills them and posts a message per group, for example
        points for an exercise sheet (\autoref{ch:groups}),
  \item saves a course layout as a template for the next semester (\autoref{ch:templates}),
  \item archives a course at the end of the semester and prepares the next one,
  \item runs your own chat bots, if you write any (\autoref{ch:bots}).
\end{itemize}

There is no web interface to learn. Everything happens in one chat, and the bot explains
each command when you ask: \cmd{!help}.

\section{Roles}
\label{sec:roles}

What you can do depends on your role. Roles are stored in the bot, not in Matrix itself.

\begin{center}
\begin{tabularx}{\textwidth}{@{}lLL@{}}
\toprule
\textbf{Role} & \textbf{Who gives it} & \textbf{What it allows} \\
\midrule
Professor & a bot admin (\cmd{!admin prof add}) & create courses and own them; appoint TAs; templates; own bots \\
Course owner & automatic: whoever creates the course & everything in that course, incl. archive and rollover \\
TA & the course owner (\cmd{!course ta add}) & in that course only: invite people, signup links, roster, groups \\
Bot admin & the server operators (bot configuration) & everything, in every course; approve bots; password resets for others; guests \\
\bottomrule
\end{tabularx}
\end{center}

A university email address alone gives no rights. Students have no role at all: the bot
only lets them manage their own password (\cmd{!password}, \autoref{sec:passwords}) and
ignores everything else they send.

\cmd{!whoami} shows your roles, and \cmd{!help} lists only the commands they allow.

\section{How to read this handbook}

Commands are shown as you type them, e.g.\ \cmd{!course info CS101}. In command
descriptions, placeholders are written in angle brackets and optional parts in square
brackets:

\begin{center}
\cmd{!course create <code> "<title>" [--semester HS26]}
\end{center}

Here you replace \cmd{<code>} and \cmd{<title>} with your own values and may leave out
\cmd{--semester HS26}. You don't type the brackets.

Conversations with the bot look like this:

\begin{you}
!course info CS101
\end{you}
\begin{botreply}
\textbf{CS101} Introduction to Computer Science (HS26)\\
Owner: @prof.example:matrix.dmi.unibas.ch\\
TAs: @ta.one:matrix.dmi.unibas.ch\\
Students (joined or invited): 84 \dots
\end{botreply}

Boxes mark tips, things to be careful about, and background notes:

\begin{tip}
\cmd{!help <command>} explains any command in the chat, with an example. You don't need this
handbook open to work with the bot.
\end{tip}

\iftechnical
\begin{note}[Technical edition]
Part~I (this part) is identical to the professors' handbook. Part~II, starting with
\autoref{ch:architecture}, describes how the bot works inside.
\end{note}
\fi
